% ============================================================
% CHAPTER 2: PROPOSED SYSTEM / ARCHITECTURE
% Trích xuất từ: 2 thesis mẫu tại thes_template/
% Cấu trúc: Related Work → Core Layers → Technology Architecture → Summary
% ============================================================

\chapter{Proposed System Architecture}
\label{ch:architecture}

% ============================================================
% SECTION 2.1: RELATED WORK AND BASE ARCHITECTURE
% ============================================================
\section{Related Work and Base Architecture}
\label{sec:related_work}

% ==========================================================
% SUBSECTION 2.1.1: Reference Architecture
% ==========================================================
\subsection{Reference Architecture}
\label{sec:reference_arch}

% PLACEHOLDER: Giới thiệu kiến trúc tham chiếu
% Tham khảo thesis 2023: vDFKM framework
% Tham khảo thesis 2026: mở rộng vDFKM

Our system architecture is referenced by a framework called [Reference Framework Name]~\citep{framework2023}. This is a standard platform based on [architecture type] applied in enterprise and organizational environments when it is necessary to build a centralized, secure, and fast data platform for in-depth analysis purposes.

\Cref{fig:reference_arch} shows the architecture of the reference framework.

\begin{figure}[tbp]
\centering
\fbox{\parbox{0.8\textwidth}{\centering\vspace{2cm}
[Insert reference architecture diagram here]\\
Architecture diagram showing layers: Data Source → Ingestion → Storage → Transformation → Metadata → Analytics
\vspace{2cm}}}
\caption{Reference framework architecture.}
\label{fig:reference_arch}
\end{figure}

% ==========================================================
% SUBSECTION 2.1.2: Proposed Architecture Overview
% ==========================================================
\subsection{Proposed High-Level Architecture}
\label{sec:proposed_arch}

% PLACEHOLDER: Mô tả kiến trúc đề xuất cao cấp
% 5 tầng chính + mô tả từng tầng

Based on the reference architecture, the proposed platform consists of five major layers that work together to support [goal]. At the base, the [Layer 1] leverages [technology] to [capability]. These changes are [processed/stored] into the platform and passed along to the next layer.

The [Layer 2] integrates [technologies] for [capability]. This layer enables both [capability 1] and [capability 2] access to raw and curated data.

Above this, the [Layer 3] utilizes [technologies] for [capability]. [Orchestration tool] orchestrates the entire pipeline, ensuring automated, reliable data workflows.

The semantic context and governance of the data are managed in the [Layer 4]. [Tool] plays a central role by [capability]. [Validation tool] enhances this layer with data validation and profiling.

At the top, the [Layer 5] enables end-users to interact with the data through [tools]. Integrated [technology] allows users to pose natural language questions, which are translated into executable queries.

\Cref{fig:proposed_arch} illustrates the proposed high-level architecture.

\begin{figure}[tbp]
\centering
\fbox{\parbox{0.9\textwidth}{\centering\vspace{2.5cm}
[Insert proposed architecture diagram here]\\
Show all 5 layers with arrows indicating data flow
\vspace{2.5cm}}}
\caption{Proposed high-level architecture.}
\label{fig:proposed_arch}
\end{figure}

% ============================================================
% SECTION 2.2: CORE LAYERS OF THE DATA PLATFORM
% ============================================================
\section{Core Layers of the Data Platform}
\label{sec:core_layers}

% ==========================================================
% 2.2.1 Data Source Layer
% ==========================================================
\subsection{Data Source Layer}
\label{sec:data_source_layer}

The Data Source Layer is where all the platform's data sources are centered. It can include a variety of different sources, ranging from relational databases to APIs, IoT sensors, web logs, or enterprise applications. These data sources can be structured, semi-structured, or unstructured.

\textbf{Structured data} is organized in a specific format (tables, columns), with predefined constraints, and easy to manage through traditional database management systems.

\textbf{Semi-structured data} has some organization but not enough to be considered structured. Formats such as XML, JSON, and HTML are common.

\textbf{Unstructured data} typically lacks any rules or organization and can be in various formats such as social media posts, image files, audio files, or unstructured text.

\Cref{tab:data_sources} summarizes the data sources used in this thesis.

\begin{table}[tbp]
\centering
\caption{Data sources used in the platform.}
\label{tab:data_sources}
\begin{tabular}{@{}lll@{}}
\toprule
\textbf{Source} & \textbf{Type} & \textbf{Format} \\
\midrule
Source A & Relational DB & Structured \\
Source B & API & Semi-structured \\
Source C & File Storage & Unstructured \\
\bottomrule
\end{tabular}
\end{table}

% ==========================================================
% 2.2.2 Data Ingestion Layer
% ==========================================================
\subsection{Data Ingestion Layer}
\label{sec:data_ingestion_layer}

The Data Ingestion Layer is the bridge between the Data Source Layer and downstream systems. This layer performs the process of importing, acquiring data from various sources for further analysis.

Data Ingestion tasks can be divided into two types:

\begin{description}
    \item[\textbf{Batch ingestion:}] Processes large volumes of data at once, typically on a scheduled basis. This method is useful when large amounts of data need to be ingested.
    \item[\textbf{Streaming ingestion:}] Processes data as it is generated, in real-time. This method is useful when data needs to be analyzed in real-time, such as in fraud detection or sensor monitoring.
\end{description}

Tools for data ingestion include open-source solutions like [Tool A], [Tool B], and commercial offerings from vendors like [Vendor A] and [Vendor B].

% ==========================================================
% 2.2.3 Data Storage and Lakehouse Layer
% ==========================================================
\subsection{Data Storage and Lakehouse Layer}
\label{sec:storage_layer}

The Data Storage and Lakehouse Layer is the foundational component of any modern data platform. This layer is responsible for storing and managing data, including raw, semi-structured, and processed datasets.

There are two popular types of storage architecture: \textbf{block storage} (used for high-performance transactional systems) and \textbf{object storage} (designed for scalable, distributed data access).

Object storage stores data in individual units at the same level. It does not use folders or subdirectories to organize data. Each object includes metadata, which provides information about the file and helps with its processing and usability.

\textbf{Data lake}, \textbf{data warehouse}, and \textbf{data lakehouse} are three different approaches to managing and processing large amounts of data:

\begin{itemize}
    \item \textbf{Data lake:} Storage repository for large amounts of raw data in any format.
    \item \textbf{Data warehouse:} Centralized repository optimized for querying and reporting.
    \item \textbf{Data lakehouse:} Hybrid architecture combining the best of both data lakes and data warehouses, enabling ACID transactions.
\end{itemize}

% ==========================================================
% 2.2.4 Data Transformation Layer
% ==========================================================
\subsection{Data Transformation Layer}
\label{sec:transformation_layer}

The Data Transformation Layer plays a critical role in the data stream by converting raw data into clean and structured formats suitable for analytics and downstream processing.

Transformation involves:

\begin{itemize}
    \item \textbf{Data cleaning:} Removing unwanted or irrelevant data, dropping unused columns, filtering duplicate values.
    \item \textbf{Data normalization:} Ensuring values are represented in consistent units and formats.
    \item \textbf{Data aggregation:} Calculating metrics such as averages, sums, and counts.
    \item \textbf{Schema alignment:} Mapping varying structures into unified formats.
\end{itemize}

Modern transformation workflows may also incorporate real-time streaming capabilities to enable continuous updates of data lakes.

% ==========================================================
% 2.2.5 Metadata and Knowledge Graph Layer
% ==========================================================
\subsection{Metadata and Knowledge Graph Layer}
\label{sec:metadata_layer}

Metadata and knowledge graph layer is one of the most important layers in data platforms by organizing and contextualizing metadata with meaningful links. This layer integrates with all the other layers to aggregate metadata throughout the system and visualize them into a knowledge graph.

This layer will recognize entities such as datasets, jobs, dashboards, and the relationships between them. Advanced properties of entities such as owners or tags will also need to be carefully configured.

From a practical standpoint, the Metadata and Knowledge Graph Layer supports advanced capabilities such as data discovery, impact analysis, and recommendation systems. It provides access to information about datasets, including schemas, data quality, update frequency, business definitions, ownership, and lineage.

\Cref{fig:kg_layer} shows the knowledge graph structure in the platform.

\begin{figure}[tbp]
\centering
\fbox{\parbox{0.7\textwidth}{\centering\vspace{2cm}
[Insert knowledge graph diagram here]\\
Nodes: datasets, jobs, dashboards; Edges: lineage, ownership, uses
\vspace{2cm}}}
\caption{Knowledge graph structure in the platform.}
\label{fig:kg_layer}
\end{figure}

% ==========================================================
% 2.2.6 Analytics and Visualization Layer
% ==========================================================
\subsection{Analytics and Visualization Layer}
\label{sec:analytics_layer}

This layer includes visualization tools that allow users to explore and comprehend their data visually. These technologies help users uncover trends, patterns, and anomalies in their data and efficiently communicate insights to stakeholders.

An important aspect in this layer supporting analytics tasks is the [LLM module], which helps translate user questions into queries and commands to generate visualizations. Instead of queries like ``SELECT SUM(...) FROM...'', users can simply ask the module to perform specific analytical tasks.

This approach makes it easier for people to get started and allows more people across the organization to access data. These approaches work best when paired with strong semantic storage that provides aggregated metrics.

% ==========================================================
% 2.2.7 Data Flow in the Platform
% ==========================================================
\subsection{Data Flow in the Platform}
\label{sec:data_flow}

\Cref{fig:data_flow} illustrates the complete data flow within the platform.

\begin{figure}[tbp]
\centering
\fbox{\parbox{0.9\textwidth}{\centering\vspace{2cm}
[Insert data flow diagram here]\\
Show: Source → Ingestion → Storage → Transformation → Analytics → Visualization
\vspace{2cm}}}
\caption{Data flow in the platform.}
\label{fig:data_flow}
\end{figure}

The data flow within the proposed platform embodies a cohesive pipeline that transforms raw, heterogeneous data into structured, semantically rich insights. It begins at the Data Source Layer, where data is ingested from multiple origins. The data then traverses the Data Ingestion Layer, which standardizes and streams the inputs. Upon ingestion, data is persisted in the Data Storage Layer. The Data Transformation Layer applies cleansing, normalization, and enrichment logic. Finally, insights propagate to the Analytics Layer where dashboards and reports are generated.

% ============================================================
% SECTION 2.3: TECHNOLOGY ARCHITECTURE
% ============================================================
\section{Technology Architecture}
\label{sec:tech_architecture}

% ==========================================================
% SUBSECTION 2.3.1: Data Source Technologies
% ==========================================================
\subsection{Data Source Technologies}
\label{sec:tech_source}

% PLACEHOLDER: MySQL, SQL Server
% Tham khảo thesis 2023/2026: mô tả chi tiết 2 DBMS

\textbf{MySQL:} MySQL is an open-source database management system. It is widely used due to its ease of deployment, speed, and compatibility across various systems. MySQL supports a wide range of storage engines and replication techniques, making it suitable for distributed environments and data-driven applications.

\textbf{Microsoft SQL Server:} Microsoft SQL Server is a relational database engine designed to retrieve, store, and manage data for structured querying, transactional processing and business intelligence features. SQL Server integrates seamlessly with enterprise tools such as Power BI and Azure Synapse.

% ==========================================================
% SUBSECTION 2.3.2: Data Ingestion Technologies
% ==========================================================
\subsection{Data Ingestion Technologies}
\label{sec:tech_ingestion}

% PLACEHOLDER: Debezium + Kafka / Airbyte
% Tham khảo thesis 2023: Airbyte
% Tham khảo thesis 2026: Debezium + Kafka

\textbf{[Tool A, e.g., Debezium or Airbyte]:} [Tool] is an open-source [type of tool] that helps [what it does]. One of the key advantages of [tool] is its capacity to manage a wide range of data sources and formats, including databases, APIs, and flat files.

\Cref{fig:ingestion_arch} shows the architecture of [tool].

\begin{figure}[tbp]
\centering
\fbox{\parbox{0.7\textwidth}{\centering\vspace{2cm}
[Insert tool architecture diagram here]
\vspace{2cm}}}
\caption{[Tool] architecture.}
\label{fig:ingestion_arch}
\end{figure}

\textbf{Apache Kafka:} Apache Kafka is a distributed event streaming platform that serves as the backbone in real-time data pipelines and streaming applications. Kafka is increasingly adopted in real-time analytics applications and event-driven architectures due to its ability to handle large volumes of data with high throughput.

Key components of Kafka include:
\begin{itemize}
    \item \textbf{Producers:} Clients that publish data to Kafka topics.
    \item \textbf{Topics:} Logical channels to which records are written.
    \item \textbf{Brokers:} Kafka servers that handle incoming data and client requests.
    \item \textbf{Consumers:} Applications that subscribe to topics and process messages.
\end{itemize}

% ==========================================================
% SUBSECTION 2.3.3: Storage and Lakehouse Technologies
% ==========================================================
\subsection{Storage and Lakehouse Technologies}
\label{sec:tech_storage}

% PLACEHOLDER: MinIO + Delta Lake / Apache Hudi

\textbf{MinIO:} MinIO is an open-source object storage system ideal for data storage applications such as data lakes and lakehouses. MinIO's distributed architecture enables it to scale horizontally, adding storage space and throughput as required.

\textbf{[Format Framework, e.g., Delta Lake or Apache Hudi]:} [Format] is a [description]. It brings database-like functionalities to large-scale data lakes, introducing capabilities such as record-level insert/update/delete (upsert), change data capture (CDC), and time travel queries.

% ==========================================================
% SUBSECTION 2.3.4: Data Transformation and Orchestration
% ==========================================================
\subsection{Data Transformation and Orchestration}
\label{sec:tech_transformation}

% PLACEHOLDER: Apache Spark + Apache Airflow

\textbf{Apache Spark:} Apache Spark is a widely used data processing engine particularly well-suited for large-scale data transformation use cases. Spark provides a fast and flexible platform for processing data in parallel across distributed clusters.

Key components of Spark:
\begin{itemize}
    \item \textbf{Driver Program:} The central coordinator that translates user code into tasks.
    \item \textbf{Cluster Manager:} Allocates resources across Spark applications.
    \item \textbf{Executors:} Distributed agents responsible for executing tasks.
\end{itemize}

\textbf{Apache Airflow:} Apache Airflow is an open-source platform for developing, scheduling, and monitoring batch workflows. It provides a Python-based framework where complex processes are expressed as Directed Acyclic Graphs (DAGs).

% ==========================================================
% SUBSECTION 2.3.5: Metadata and Knowledge Graph Management
% ==========================================================
\subsection{Metadata and Knowledge Graph Management}
\label{sec:tech_metadata}

% PLACEHOLDER: Datahub + Great Expectations

\textbf{DataHub:} DataHub is a modern data catalog designed to streamline metadata management, data discovery, and data governance. It enables efficient exploration of organizational data, lineage tracking, and enforcement of data contracts. DataHub models metadata as a graph of interconnected entities (datasets, dashboards, users) and relationships (schema, ownership, lineage).

\textbf{Great Expectations:} Great Expectations is an open-source data validation, documentation, and profiling tool designed to help teams maintain data quality throughout the data pipeline. It offers a declarative framework where users define expectations---assertions about data characteristics---that are automatically checked during data ingestion, transformation, and serving.

% ==========================================================
% SUBSECTION 2.3.6: Querying and Analytics
% ==========================================================
\subsection{Querying and Analytics}
\label{sec:tech_querying}

% PLACEHOLDER: Trino + Metabase / Superset

\textbf{Trino:} Trino, formerly known as PrestoSQL, is an open-source distributed SQL query engine designed for running fast, interactive analytic queries across heterogeneous data sources. It allows querying data directly where it lives---whether on object storage systems like S3, distributed file systems like HDFS, or in traditional relational databases---without requiring prior data movement.

The Trino architecture follows a distributed design based on a coordinator-worker model:
\begin{itemize}
    \item \textbf{Coordinator:} Parses SQL statements, optimizes query plans, schedules tasks.
    \item \textbf{Workers:} Execute parts of the query in parallel.
    \item \textbf{Catalogs:} Provide mappings to external data sources.
\end{itemize}

\textbf{[BI Tool, e.g., Metabase or Superset]:} [Tool] is an open-source business intelligence and data visualization platform designed to make data exploration accessible to non-technical users through an intuitive, low-code interface.

% ==========================================================
% SUBSECTION 2.3.7: LLM-Powered Analytics (Optional)
% ==========================================================
\subsection{LLM-Powered Intelligent Analytics (Optional)}
\label{sec:tech_llm}

The emergence of Large Language Models (LLMs) has fundamentally transformed the way users interact with data systems, introducing natural language interfaces that enable intuitive and dynamic data exploration.

In the context of intelligent analytics platforms, LLMs play a pivotal role by interpreting user prompts, automatically generating SQL queries, recommending visualizations, and providing narrative summaries of analytical results.

The architecture for LLM-powered intelligent analytics integrates LLMs into the data platform via a specialized interaction layer:
\begin{itemize}
    \item \textbf{Large Language Models:} Serve as the core natural language understanding and generation engine.
    \item \textbf{Integration Layer:} Acts as a bridge between the LLM and backend services.
\end{itemize}

% ============================================================
% SECTION 2.4: CHAPTER SUMMARY
% ============================================================
\section{Chapter Summary}
\label{sec:ch2_summary}

This chapter introduced the core layers and technology stack underpinning the proposed [Knowledge Graph-based / Data Fabric-based] Data Platform for [Intelligent Analytics / Modern Data Management]. We explored each component's role from data storage and transformation, to metadata management, intelligent querying, and LLM-powered analytics.

The proposed architecture consists of [5-7] major layers that work together to provide [capabilities]. The technology stack includes open-source tools such as [list of tools], which were selected based on their [criteria: scalability, flexibility, integration capability, community support].

Together, these technologies form an integrated ecosystem that enables efficient, scalable, and intelligent data workflows, providing a strong foundation for building advanced analytics capabilities within modern data-driven organizations. The following chapter will present the experimental setup and evaluation results.
